\documentclass[../main.tex]{subfiles}

\begin{document}

\section{Resources and Supporting Tooling}
\label{app:resources}

The artifacts below are referenced by recommendations in this report as
existing starting points. Listing a project here identifies a resource; it is
not an assessment of that project, and none was audited.

\begin{description}[leftmargin=1.5em,style=nextline]
    \item[HDF5 Safety, Security \& Privacy (SSP) Special Interest
    Group~\cite{hdf5-ssp-sig}.]
    Governance, process, and artifact home for the work this report feeds.
    It is the natural venue for the roadmap adoption decision described in
    Appendix~\ref{app:phased-roadmap}, and for publishing the
    \code{CORE-S12} metrics dashboard.

    \item[HDF5 CVE test suite~\cite{cve_hdf5}.]
    An existing collection of vulnerability-triggering inputs. It is the
    obvious seed for the permanent corpus proposed in \code{CORE-S6}, which
    additionally requires per-case expected-outcome metadata and execution in
    both assertion-enabled and \code{-DNDEBUG} builds. It also contributed to
    the CVE enumeration in Appendix~\ref{app:cve-history-summary}.

    \item[\code{hdf5-pickles}: a machine-readable specification of the HDF5
    file format~\cite{hdf5-pickles}.]
    Directly relevant to three recommendations. It is a candidate source of
    truth for the invariant registry in \code{CORE-S1}; a generator for the
    positive and negative fixtures required by \code{IND-S1}; and the origin
    of issue~\#87, the assert-masking observation tracked as
    \code{CORE-SEC-003}.

    \item[HDF5 extension skeletons~\cite{hdf5-boneyard}.]
    Templates for VOL connectors, filters, and VFDs. A natural place to
    embed the privacy-documentation template required by \code{FILTER-S5} and
    the trust and search-path defaults required by \code{VOL-S3},
    \code{FILTER-S1}, and \code{VFD-S4}, so that new extensions inherit them
    rather than retrofitting them.

    \item[HDF5 plugins~\cite{hdf5_plugins}.]
    The filter-plugin project whose options govern the dependency set
    discussed in Section~\ref{sec:supply-chain} and
    Appendix~\ref{app:supply-chain}. It is the scope boundary for the plugin
    manifest and allowlist proposed in \code{CORE-S10}.
\end{description}

\end{document}
